\chapter{Penjumlahan: Langkah Pertama}\label{ch:g1:addition}

Tiga kelereng merah dan dua kelereng biru: berapa banyak semuanya?
Menggabungkan lalu membilang seluruhnya berarti \emph{menjumlahkan}
--- operasi yang pertama, dengan dua tandanya yang terkenal $+$ dan
$=$.

\section{Menggabungkan}

\begin{definition}[Penjumlahan]\label{def:g1:addition:def}
\emph{Menjumlahkan} berarti menggabungkan lalu membilang semuanya.
Kita menulis
\[
3 + 2 = 5
\]
dan dibaca: ``tiga tambah dua sama dengan lima''. Hasil sebuah
penjumlahan disebut \emph{jumlah}\index{jumlah}.
\end{definition}

\begin{omfigure}
\begin{tikzpicture}[scale=0.6]
  \foreach \i in {0,1,2} \fill[omThm] (\i*1.0,0) circle (0.3);
  \node[font=\large] at (3.0,0) {$+$};
  \foreach \i in {0,1} \fill[omDef] (3.8+\i*1.0,0) circle (0.3);
  \node[font=\large] at (5.8,0) {$=$};
  \foreach \i in {0,1,2} \fill[omThm] (6.6+\i*1.0,0) circle (0.3);
  \foreach \i in {0,1} \fill[omDef] (9.6+\i*1.0,0) circle (0.3);
  \node[below, font=\small] at (5.6,-0.7)
    {$3 + 2 = 5$: tiga dan dua menjadi lima};
\end{tikzpicture}

{\small Penjumlahan dalam gambar: kedua kelompok bersama-sama menjadi
satu kelompok berisi lima.}
\end{omfigure}

\begin{example}[Urutannya tidak berpengaruh]\label{ex:g1:addition:order}
$2 + 3$ dan $3 + 2$ menggabungkan kelereng yang sama: keduanya
menjadi $5$. Untuk menghitung $2 + 9$, lebih cepat memulai dari
bilangan yang lebih besar: sebutlah ``sembilan'', lalu bilang maju
dua --- ``sepuluh, sebelas''. Jadi $2 + 9 = 11$.
\end{example}

\section{Sahabat sepuluh}

\begin{example}[Pasangan yang berjumlah 10]\label{ex:g1:addition:friends}
Dua bilangan yang bersama-sama menjadi $10$ disebut ``sahabat sepuluh'':
\[
1 + 9, \quad 2 + 8, \quad 3 + 7, \quad 4 + 6, \quad 5 + 5 .
\]
Sepuluh jari memperlihatkan semuanya: angkatlah $3$ jari, maka $7$
jari yang terlipat adalah sahabatnya. Hafal di luar kepala sahabat
sepuluh membuat banyak penjumlahan menjadi cepat.
\end{example}

\begin{omfigure}
\begin{tikzpicture}[scale=0.72]
  % bingkai sepuluh: 2 baris berisi 5 kotak
  \draw[thick] (0,0) grid[step=1] (5,2);
  \draw[very thick] (0,0) rectangle (5,2);
  \foreach \i in {0,1,2} \fill[omDef] (\i+0.5,1.5) circle (0.3);
  \foreach \i in {3,4} \draw[omThm, thick] (\i+0.5,1.5) circle (0.3);
  \foreach \i in {0,...,4} \draw[omThm, thick] (\i+0.5,0.5) circle (0.3);
  \node[right, font=\small, align=left] at (5.6,1.0)
    {sebuah \emph{bingkai sepuluh}: $10$ kotak.\\
     $3$ terisi, $7$ kosong: $3 + 7 = 10$};
\end{tikzpicture}

{\small Bingkai sepuluh memperlihatkan setiap pasangan sahabat sepuluh
sekaligus: yang terisi dan yang kosong selalu berjumlah $10$.}
\end{omfigure}

\begin{method}[Melewati sepuluh]\label{met:g1:addition:crossten}
Untuk menghitung $8 + 5$:
\begin{enumerate}
  \item ambillah dari $5$ sebanyak yang dibutuhkan $8$ untuk mencapai
        $10$, yaitu $2$ (sahabat sepuluhnya);
  \item $8 + 2 = 10$, dan tersisa $3$ dari $5$ tadi;
  \item $10 + 3 = 13$. Jadi $8 + 5 = 13$.
\end{enumerate}
Melewati $10$ mengubah satu \omterm{def:g1:addition:def}{jumlah} yang sulit menjadi dua yang mudah.
\end{method}

\begin{omfigure}
\begin{tikzpicture}[scale=0.62]
  % bingkai sepuluh dengan 8 biru, 2 merah (dipinjam dari yang 5)
  \draw[thick] (0,0) grid[step=1] (5,2);
  \draw[very thick] (0,0) rectangle (5,2);
  \foreach \i in {0,...,4} \fill[omDef] (\i+0.5,1.5) circle (0.3);
  \foreach \i in {0,1,2} \fill[omDef] (\i+0.5,0.5) circle (0.3);
  \foreach \i in {3,4} \fill[omThm] (\i+0.5,0.5) circle (0.3);
  % 3 merah tersisa di luar bingkai
  \foreach \i in {0,1,2} \fill[omThm] (6.2+\i,1.0) circle (0.3);
  \node[below, font=\small] at (2.5,-0.25) {bingkainya penuh: $10$};
  \node[below, font=\small] at (7.2,0.55) {$3$ meluber};
\end{tikzpicture}

{\small $8 + 5$ pada bingkai sepuluh: $8$ keping biru mengambil $2$
dari $5$ keping merah untuk memenuhi bingkai --- satu puluhan penuh
dan $3$ lagi, jadi $13$.}
\end{omfigure}

\begin{omfigure}
\begin{tikzpicture}[scale=0.6]
  \draw[-Stealth] (-0.4,0) -- (15.4,0);
  \foreach \i in {0,...,15} \draw (\i,-0.1) -- (\i,0.1);
  \foreach \i in {0,5,8,10,13,15}
    \node[below, font=\footnotesize] at (\i,-0.12) {\i};
  \draw[-Stealth, omDef, thick] (8,0.3) .. controls (9,0.95) ..
    (10,0.3);
  \node[omDef, font=\small] at (9,1.05) {$+2$};
  \draw[-Stealth, omThm, thick] (10,0.3) .. controls (11.5,1.15) ..
    (13,0.3);
  \node[omThm, font=\small] at (11.5,1.2) {$+3$};
\end{tikzpicture}

{\small $8 + 5$ pada garis bilangan: melompat ke $10$ dulu, lalu sisanya.}
\end{omfigure}

\begin{example}[Jumlah kembar]\label{ex:g1:addition:doubles}
\omterm{def:g1:addition:def}{Jumlah} kembar layak dihafal di luar kepala:
\[
1+1=2, \quad 2+2=4, \quad 3+3=6, \quad 4+4=8, \quad 5+5=10,
\]
\[
6+6=12, \quad 7+7=14, \quad 8+8=16, \quad 9+9=18, \quad 10+10=20 .
\]
Yang hampir kembar didapat cuma-cuma: $6 + 7$ satu lebihnya dari $6 + 6$, jadi $13$.
\end{example}

\section{Soal cerita}

\begin{example}\label{ex:g1:addition:problem}
``Lia punya $7$ stiker. Ia mendapat $6$ lagi. Berapa sekarang?''
Mendapat tambahan berarti menjumlahkan: $7 + 6 = 13$ (lewat sepuluh:
$7 + 3 = 10$, lalu $+3$). Lia punya $13$ stiker. Selalu akhiri dengan
satu kalimat jawaban!
\end{example}

\section{Latihan}

\begin{exercise}[$\star$]\label{exo:g1:addition:1}
Gambarlah kelerengnya lalu hitunglah: $4 + 3$; \quad $5 + 2$; \quad
$6 + 4$.
\end{exercise}

\begin{exercise}[$\star$]\label{exo:g1:addition:2}
Hitunglah dengan memulai dari bilangan yang lebih besar: $3 + 8$;
\quad $2 + 12$; \quad $4 + 9$.
\end{exercise}

\begin{exercise}[$\star$]\label{exo:g1:addition:3}
Sebutkan sahabat sepuluh dari: $7$; \ $4$; \ $9$; \ $5$; \ $1$.
\end{exercise}

\begin{exercise}[$\star$]\label{exo:g1:addition:4}
Hitunglah dengan melewati sepuluh (\cref{met:g1:addition:crossten}):
$9 + 4$; \quad $8 + 7$; \quad $6 + 5$.
\end{exercise}

\begin{exercise}[$\star$]\label{exo:g1:addition:5}
Sebutkan \omterm{def:g1:addition:def}{jumlah} kembar, lalu pakailah untuk menghitung: $7 + 8$; \quad
$5 + 6$; \quad $9 + 8$.
\end{exercise}

\begin{exercise}[$\star$]\label{exo:g1:addition:6}
Lengkapilah bagian yang kosong:
\[
6 + \;?\; = 10, \qquad
\;?\; + 5 = 12, \qquad
10 + \;?\; = 17 .
\]
\end{exercise}

\begin{exercise}[$\star$]\label{exo:g1:addition:7}
``Di dalam bus ada $8$ anak; $5$ orang dewasa naik. Berapa orang yang
ada di dalam bus?'' Tulislah penjumlahannya dan kalimat jawabannya.
\end{exercise}

\begin{exercise}[$\star$]\label{exo:g1:addition:8}
Tono punya $6$ mobil merah dan $7$ mobil biru. Mia punya $10$ mobil.
Siapa yang mobilnya lebih banyak?
\end{exercise}

\begin{exercise}[$\star\star$]\label{exo:g1:addition:9}
Tiga dadu menunjukkan $4$, $6$ dan $5$. Berapa \omterm{def:g1:addition:def}{jumlahnya}? (Jumlahkan
dua dadu dulu --- pilihlah dua yang menjadi sahabat sepuluh!)

\end{exercise}
